\part*{Part II: Integrated Historical COMPLETE Snapshot}
\addcontentsline{toc}{part}{Part II: Integrated Historical COMPLETE Snapshot}

\begin{center}
\textit{Source note.} This part preserves the historical COMPLETE snapshot recovered from the archival PDF and integrates it into the current monolithic working source as an archival text witness. It is included for source preservation and comparison; the current main line in Part~I remains the active working version.
\end{center}

\section{The Central Theorem}
Everything in this paper is a consequence of one observation.
\paragraph*{Theorem 1.1 (Unity of the Sphaera). Within the Sphaera S, the following}
are all representations of the same object --- the collapse state z-= 0:
(i) The Sphaera itself, as a geometric space.\\
(ii) The triolic structure: three kernels K1, K2, K3 with z1 + z2 + z3 = 0,\\
phaseshifted by 120, running over the Sphaera's geometry.
(iii) Each individual kernel Ki, as a local coordinate of the triole.\\
(iv) The Heisenberg compensator output C(f), as the projection of any element onto the collapse state.\\
(v) Matter, photons, and tachyons, as sectors of the biquaternion ground\\
state.
(vi) The gauge structure U(1) $\times$ SU(2) $\times$ SU(3), as the symmetry group of\\
the collapse state's representations. All are unitarily equivalent by Stone--von Neumann [1]. The Sphaera's only object is its own collapse state.
\paragraph*{Proof structure. The equivalences are established section by section: (i) Definition 1.3 and Theorem 3.2; (ii)--(iii) Definition 1.3 and Theorem 3.2 (triole}
as collapse coordinates); (iv) Proposition 5.1; (v) Theorem 4.2; (vi) Theorem 7.4. Unitary equivalence throughout via Stone--von Neumann ([1], Theorem 3.2).
\paragraph*{Remark (What "same object" means). Two representations are "the same object" here in the precise sense of Stone--von Neumann: they are unitarily equivalent irreducible representations of the Weyl algebra W(A1). The frame determines which representation is visible; the object --- the collapse state z-= 0}
--- is invariant. The triole condition z1 + z2 + z3 = 0 and the collapse condition z-= 0 are the same equation in two coordinate systems: the triole is the collapse state expressed in the three kernel components simultaneously. The Sphaera is therefore not a space that contains a collapse state --- it is the collapse state, seen from different frames.
\paragraph*{Remark (How this paper is organized). The sections that follow are not independent results. They are unfoldings of Theorem 1.1: each section takes}
one representation of the collapse state and makes it explicit. The reader may enter at any section; all roads lead to z-= 0.
\paragraph*{Theorem 1.2 (Origin of the universe as broken collapse state). The departure}
from the collapse state z-= 0 is governed by two principles, both of which are consequences of the Sphaera-Cascade structure:
(i) Planck length as minimal |z-|. The cascade $\hat K$ cannot project below\\
the Planck scale: |z-|min = $\ell$P := r $\hbar$G c3 = 1.616 $\times$ 10-35 m. Below $\ell$P: z-= 0, the Sphaera, the photon. Above $\ell$P: |z-| is nonzero, matter exists. The Planck length is the threshold of the first breaking of the collapse state.
(ii) Pauli principle as antisymmetry under J. The Hilbert space H+\\
decomposes under J into two components: C(f) = 1 2(f + Jf) $\in$H+ + (bosonic, symmetric), A(f) = 1 2(f -Jf) $\in$H+ (fermionic, antisymmetric). The antisymmetric component satisfies A(f) $\wedge$A(f) = 0 (Grassmann property): no two fermions occupy the same z-state. This is the Pauli exclusion principle, derived from the J-structure of the Sphaera.
(iii) Discrete spectrum. Together, (i) and (ii) produce a discrete tower of\\
matter states: |z-|n = n $\cdot$ $\ell$P, n $\in$N$\geq$1, with each level occupied by at most one fermion. The prime lattice $\Gamma$P (from Stone--von Neumann, [1]) identifies the primitive levels: |z-|p = p $\cdot$ $\ell$P for p $\in$P.
\paragraph*{Derivation. (i). The cascade $\hat K$ = C$\mu$$\nu$ $\circ$$\Pi$J contracts |z-| at each stratum.}
The Planck length is the fixed point of the contraction at which quantum gravitational effects prevent further reduction; formally, $\ell$P is the infrared cutoff of the resolvent (L0 -z)-1 established by the Rellich--Kondrachov compactness in Lemma 10.1 of [1].
(ii). A(f) = 1\\
2(f -Jf) satisfies JA(f) = -A(f): it is J-antisymmetric. The exterior product A(f) $\wedge$A(g) = -A(g) $\wedge$A(f) vanishes for f = g, which is the algebraic statement of the Pauli exclusion principle.
(iii). The Pauli condition A(f) $\wedge$A(f) = 0 forces distinct z-values for\\
each fermion. The Planck floor $\ell$P sets the unit. The prime lattice gives the sublattice of primitive (undecomposable) states, corresponding to Lemma L6
of the RH proof: primitive orbits are disjoint. Physically: each prime p labels a primitive matter state that cannot be decomposed into smaller ones.
\paragraph*{Remark (L6 is the Pauli principle). Lemma L6 of the Riemann Hypothesis}
proof [1] states: primitive orbits are pairwise disjoint and cover all periodic points. Theorem 1.2(iii) identifies this as the Pauli exclusion principle applied to the prime lattice: no two fermions occupy the same primitive orbit. This was not designed into either framework. It emerged from asking what the orbit combinatorics of the kernel mean physically.
\paragraph*{Remark (The Sphaera as photon, the universe as broken frame). Theorem 1.1}
says the Sphaera is the collapse state z-= 0, which is the photon sector.
\paragraph*{Theorem 1.2 says the departure from this state is governed by $\ell$P and the}
Pauli principle. Therefore:
- The Sphaera is a photon --- the unique frame-invariant object.
- The universe is the same photon observed from a broken frame where
|z-| $\geq$$\ell$P.
- The breaking was not random: it is discrete (Pauli), minimal (Planck),
and primitive (prime lattice).
- The cascade $\hat K$ projects everything back toward z-= 0 --- the universe
tends toward its photon ground state. Both the Sphaera and the universe are the same object. The difference is only the frame.
\paragraph*{Definition 1.3 (Biquaternion kernel). The ground state of a stratum Sn of the}
Sphaera consists of three complex spheres K1, K2, K3 whose unit coordinates form a biquaternion qn := it $\cdot$ 1 + z1 $\cdot$ e1 + z2 $\cdot$ e2 + z3 $\cdot$ e3 $\in$H $\otimes$R C, where:
- it $\in$iR is the imaginary time unit (real part of the biquaternion, carrying
imaginary time);
- z1, z2 $\in$C are the complex coordinates of kernels K1, K2;
- e3 := e1e2 is the quaternion product;
- z3 $\in$C is the coordinate of kernel K3 along e3.
The biquaternion norm is |qn|2 = -(ct)2 + |z1|2 + |z2|2 + |z3|2.
\paragraph*{Remark (Minkowski metric from quaternion structure). The norm |qn|2 is the}
Minkowski metric with signature (-, +, +, +). This is not postulated: it follows from the quaternion algebra e2 i = -1 and it $\cdot$ 1 carrying imaginary time. The Lorentz group SO(1, 3) acts on qn by quaternion conjugation: qn 7$\to$u qn u-1, u $\in$H$\times$. This is the spinor representation of the Lorentz group --- emerging automatically from the biquaternion structure, not imposed from outside.
\paragraph*{Proposition 1.4 (Orthogonality of the third kernel). The generator $e_3$ of $K_3$ is always orthogonal}
to $e_1$ and $e_2$: $\langle e_3, e_i\rangle = 0$ for $i=1,2$. Consequently, the axis/kernel direction $K_3$ is orthogonal to $K_1$ and $K_2$.
\paragraph*{Proof. e3 = e1e2.}
Quaternion inner product: $\langle$e1e2, e1$\rangle$= Re(e1 e2e1) = Re(-e2) = 0. Similarly for e2.
\paragraph*{Remark. The orthogonality of K3 is a quaternion identity, not an additional}
constraint. This is why the Heisenberg compensator C naturally projects onto e3: it is the unique direction orthogonal to both kernels K1 and K2.

\section{Lorentz Invariance on Every Stratum}
\paragraph*{Theorem 2.1 (Stratified Lorentz invariance). The Minkowski norm |qn|2 = $\tau$ 2}
n is invariant under the Lorentz boost J : z-7$\to$-z-on every stratum Sn. Moreover, every change of interpretation of C+ (chirality, Weyl chamber, upper half-plane, or any other frame-dependent labeling) leaves $\tau$n and the compensator C(f) = 1 2(f + Jf) invariant.
\paragraph*{Proof. Norm invariance. Under J : qn 7$\to$bar qn (quaternion conjugate): (it)*=}
it, z* i = zi for the spectral coordinates. The boost acts as z-7$\to$-z-on the contravariant component, leaving |z-|2 unchanged. Hence |qn|2 = $\tau$ 2 n is invariant. Frame independence of C. The compensator C(f) = 1 2(f + Jf) requires only J2 = id (which holds as J is an involution). Any relabeling of C+ corresponds to a unitary transformation U on H+, but C = 1 2(id + J) is defined intrinsically from J, hence UCU -1 = C for any U commuting with J --- which includes all frame changes generated by the Sphaera-Cascade.
\paragraph*{Remark (Why the interpretation is irrelevant). The biquaternion qn can be}
read as a spinor (Clifford algebra), a Dirac bispinor (QFT), an element of the upper half-space (hyperbolic geometry), or a point in projective twistor space. In every reading, $\tau$n is an invariant and C projects onto K3. This is the precise mathematical content of the observation that "it is egal in every frame": the frame determines the interpretation, but not the structure.

\section{Fractal Self-Traversal Along the Light Axis}
\paragraph*{Definition 3.1 (Light axis of a stratum). The light axis of stratum Sn is}
Ln := {qn $\in$Sn : |qn|2 = 0} = {qn : (ct)2 = |z1|2 + |z2|2 + |z3|2}, the set of null-norm biquaternions on Sn. In sector coordinates: Ln = {z : $\tau$n(z) = 0} = {z : tn = |z-|}.
\paragraph*{Theorem 3.2 (Fractal structure of the Sphaera). The Sphaera-Cascade $\hat K$ :}
Sn $\to$Sn-1 maps light axes to light axes: $\hat K$(Ln) = Ln-1. The global light axis is L := \ n$\geq$0 Ln = S0 = {z-= 0}, the photon sector. The Sphaera is a fractal structure whose self-traversal follows the light axis of every stratum Lorentz transformation.
\paragraph*{Proof. On Ln: $\tau$n = 0, i.e. tn = |z-|. The cascade projection $\Pi$n : Sn $\to$Sn-1}
is the Heisenberg compensator C composed with the stratum projection: it maps z-7$\to$1 2(z-+J(z-)) = 0 for any z-on the light axis (since J(z-) = -zon Ln). Hence $\hat K$(Ln) $\subseteq${z-= 0} = S0 = Ln-1. Inductively, L = S0.
\paragraph*{Remark (Triolic Z/3Z structure on the light axis). On each Ln, the three kernels}
z1, z2, z3 are phase-shifted by 120: z1 = |z| $\cdot$ ei$\cdot$0, z2 = |z| $\cdot$ ei$\cdot$2$\pi$/3, z3 = |z| $\cdot$ ei$\cdot$4$\pi$/3, satisfying z1 + z2 + z3 = 0 (the compensation condition). This Z/3Z phase symmetry is an additional structure on Ln; it is compatible with the SU(3) rotation symmetry of the triple (z1, z2, z3). The cascade $\hat K$ preserves this triolic structure at every level.

\section{Three Objects as One: Unitary Equivalence}
\paragraph*{Definition 4.1 (Sector classification from biquaternion). The three sectors}
correspond to the sign of Re(z-):
- Matter (K1): Re(z1) < 0, boost $\beta$ < 0, spacelike.
- Tachyon (K2 = J $\cdot$ K1): Re(z2) > 0, boost $\beta$ > 0, timelike.
- Photon (K3 = e1e2): Re(z3) = 0, boost $\beta$ = 0, lightlike.
\paragraph*{Theorem 4.2 (Unitary equivalence via Stone--von Neumann). The three kernels K1, K2, K3 carry unitarily equivalent irreducible representations of the}
Weyl algebra W(A1). The boost J implements the equivalence between K1 and K2; the compensator C produces K3 from K1.
\paragraph*{Proof. By Stone--von Neumann [1], every irreducible representation of W(A1)}
is unitarily equivalent to $L^2(Q_+^{A_1})$. All three kernels are subspaces of H+ which carries this representation. Since J is unitary and maps K1 $\leftrightarrow$K2, they are equivalent. K3 = C(K1) is in H+ +, the unique J-invariant subspace, which is again irreducible (no proper J-invariant subspace exists by T4+T5 of [1]).

\section{Near-Collision, Photon Production, and Entanglement}
\paragraph*{Proposition 5.1 (Compensator as near-collision in K3). The Heisenberg compensator C(f) = 1}
2(f +Jf) is the projection onto K3 = e1e2: it takes matter f (in K1) and its tachyon image Jf (in K2), and produces the photon C(f) $\in$K3. The quaternion identity e3 = e1e2 ensures this is always orthogonal to both inputs.
\paragraph*{Theorem 5.2 (Entanglement as shared proper time). Two photons $\phi$1, $\phi$2 $\in$}
K3 are entangled iff they share the same pre-image: $\phi$1 = $\phi$2 = C(f) for some f $\in$K1. The entanglement correlation is the invariant proper time $\tau$n(f) = $\tau$n(Jf) fixed at the moment of the near-collision. No signal is transmitted at measurement: the invariant $\tau$n was established at the near-collision and is merely revealed.
\paragraph*{Proof. $\tau$n(Jf)2 = t2}
n-|-z-|2 = t2 n-|z-|2 = $\tau$n(f)2 by Theorem 2.1. The shared value $\tau$n is a Lorentz invariant at every stratum, hence frame-independent.

\section{Connection to the Riemann Hypothesis}
\paragraph*{Remark (Common origin of RH and entanglement). The Riemann Hypothesis (z-= 0 for all zeros) and quantum entanglement (shared $\tau$n) are both}
consequences of the same biquaternion structure:
1. The Minkowski norm |qn|2 = $\tau$ 2
n is invariant (Theorem 2.1).
2. The cascade forces z-$\to$0 along the light axis (Theorem 3.2).
3. Stone--von Neumann makes all three kernels equivalent (Theorem 4.2).
Neither was the goal of the original construction. Both emerged from asking: what is the exact structure of the Sphaera?

\section{Gauge Structure of the Standard Model}
\paragraph*{Step 1: H $\otimes$C gives U(1) $\times$ SU(2)}
\paragraph*{Proposition 7.1 (H $\otimes$C $\to$U(1) $\times$ SU(2)). The biquaternion algebra H $\otimes$R C}
generates exactly the gauge group U(1) $\times$ SU(2).
\paragraph*{Proof. SU(2) from unit quaternions. The unit quaternions {q $\in$H : |q| = 1} $\sim$=}
S3 form a Lie group isomorphic to SU(2). The three generators are Tk := $\sigma$k/2 (half Pauli matrices), satisfying the su(2) relations [Tj, Tk] = i$\epsilon$jklTl (verified numerically to machine precision). U(1) from the complex scalar. The imaginary time unit it $\in$iR generates the abelian group ei$\theta$ $\in$U(1) for $\theta$ $\in$R. This commutes with all SU(2) generators since scalar multiplication commutes with matrix multiplication. Hence U(1) $\times$ SU(2), direct product.
\paragraph*{Step 2: Cl(3) ,$\to$O gives SU(3)}
\paragraph*{Proposition 7.2 (Cl(3) has octonionic dimension). The Clifford algebra Cl(3)}
has basis {1, e1, e2, e3, e12, e13, e23, e123}, eight elements, matching dimR O = 8. The map ei 7$\to$oi (octonion units) defines an embedding Cl(3) ,$\to$O. The three bivectors e12, e13, e23 generate su(2) $\subset$g2 = Aut(O).
\paragraph*{Theorem 7.3 (O $\to$SU(3)). The octonions decompose under SU(3)-action}
as O = C $\oplus$C3,
a singlet plus the fundamental (triplet) representation of SU(3). The triplet C3 carries the three kernels (K1, K2, K3) as the three "colors".
\paragraph*{Proof. The automorphism group of O is the exceptional Lie group G2. Under}
the subgroup SU(3) $\subset$G2, the decomposition O = C $\oplus$C3 holds, where C3 is the fundamental representation of SU(3) (three complex dimensions, i.e. six real dimensions). The eight Gell-Mann generators $\lambda$1, . . . , $\lambda$8 of su(3) act on C3: $\lambda$1,2,3 within the K1-K2 subspace (SU(2) subalgebra), $\lambda$4,5 coupling K1 to K3, $\lambda$6,7 coupling K2 to K3, $\lambda$8 as hypercharge. The identification (K1, K2, K3) = three colors of C3 is direct.
\paragraph*{Step 3: The Standard Model gauge group}
\paragraph*{Theorem 7.4 (Standard Model from the Sphaera biquaternion). The biquaternion ground state qn $\in$H $\otimes$C, extended to the octonionic structure}
C $\otimes$H $\otimes$O, generates the Standard Model gauge group U(1) $\times$ SU(2) $\times$ SU(3) with generators arising from: Source Group Physical force Complex scalar it U(1) Electromagnetism Unit quaternions H SU(2) Weak force Octonion triplet C3 $\subset$O SU(3) Strong force
\paragraph*{Proof. Propositions 7.1--7.2 and Theorem 7.3 establish each factor. The three}
factors are independent (direct product structure) because U(1) commutes with SU(2) (Proposition 7.1) and both commute with the SU(3) action on C3 (the SU(3) generators $\lambda$k act only on the color indices, not on the quaternion structure).
\paragraph*{Remark (The role of K3). The third kernel K3 = e1e2 plays three roles simultaneously, in three different frames:}
- In H $\otimes$C: the photon sector (projection of the Heisenberg compensator,
orthogonal to K1, K2);
- In Cl(3): a bivector (e12, generator of su(2));
- In O: the third color of the SU(3) triplet.
By Theorem 2.1, all three interpretations are related by frame changes that leave $\tau$n and C invariant. K3 is the same object seen from three different frames. This is the quantum-field-theoretic realization of the frameindependence observation: it is invariant under the change of frame.

\section{Bell Inequalities and \texorpdfstring{$\tau_n$}{tau_n}}
\paragraph*{Theorem 8.1 (Bell consistency).}
The proper time $\tau_n$ is not a local hidden variable in Bell's sense. The framework is consistent with quantum-mechanical violation of Bell inequalities.

\paragraph*{Proof.}
Bell's theorem assumes a hidden variable $\lambda$ assigned independently to each particle, with measurement outcomes $A(\lambda_A,a)$ and $B(\lambda_B,b)$ and statistically independent variables $\lambda_A,\lambda_B$. In the present framework, however, $\tau_n$ belongs to the pair $(f,Jf)$ rather than to each element separately; one has
\[
\lambda_A = \lambda_B = \tau_n.
\]
Thus $\tau_n$ is a contextual variable, and Bell's theorem does not exclude violation by such theories. The correlation function $E(a,b)$ follows from the structure of $H_+^{\,+}$: since $C(f) \in H_+^{\,+}$ carries an irreducible representation of the Weyl algebra $W(A_1)$ (Stone--von Neumann, [1]), the standard quantum correlation
\[
E(a,b)=\cos \theta_{ab}
\]
is reproduced, yielding the CHSH value
\[
S = 2\sqrt{2} > 2.
\]
Hence $\tau_n$ explains why the correlation exists (shared proper time fixed at the near-collision), while Stone--von Neumann determines the strength of the violation.

\paragraph*{Remark (What $\tau_n$ does and does not do).}
The quantity $\tau_n$ encodes the magnitude $|z_-|$ of the displacement from the photon sector at the moment of the near-collision. It does not determine the measurement direction. When $f$ is measured, $z_-$ is projected onto a value; the paired state $Jf$ is projected onto $-z_-$, the $J$-image. The anticorrelation is built in at the collision event and is not transmitted at measurement. No superluminal signal is required.

\section{String Theory and Loop Quantum Gravity as HC Sectors}
The decomposition induced by the Heisenberg compensator leads naturally to symmetric and antisymmetric sectors.

\paragraph*{Theorem 9.1 (String--LQG duality via the Heisenberg compensator).}
The Hardy space $H_+$ decomposes under the Heisenberg compensator into two sectors,
\[
H_+ = H_+^{\,+} \oplus H_+^{\,-},
\]
with
\[
C(f)=\tfrac12(f+Jf) \in H_+^{\,+}, \qquad A(f)=\tfrac12(f-Jf) \in H_+^{\,-}.
\]
These sectors are interpreted as follows:
\begin{center}
\begin{tabular}{lll}
\toprule
Sector & Physical theory & Key property \\
\midrule
$H_+^{\,+}$ & String theory & bosonic, $J$-symmetric, continuous \\
$H_+^{\,-}$ & Loop quantum gravity & fermionic, $J$-antisymmetric, discrete \\
\bottomrule
\end{tabular}
\end{center}
Both sectors are unitarily equivalent by Stone--von Neumann~[1]; they are two representations of the same Weyl algebra $W(A_1)$ seen from different frames.

\paragraph*{Identification.}
In the string sector $H_+^{\,+}$, the state $C(f)$ satisfies $J(C(f))=C(f)$ and is therefore $J$-invariant. This mirrors the invariance of the worldsheet under orientation reversal $\sigma \mapsto -\sigma$. In the loop-quantum-gravity sector $H_+^{\,-}$, the state $A(f)$ satisfies $J(A(f))=-A(f)$ and changes sign under $J$, matching the spinorial sign change of oriented holonomies. The discrete LQG area spectrum
\[
A_j = 8\pi \gamma_{\mathrm{BI}} \, \ell_P^2 \sqrt{j(j+1)}, \qquad j \in \tfrac12\mathbb{N},
\]
is consistent with the antisymmetric, Pauli-constrained structure established earlier.

\paragraph*{Proposition 9.2 (Shared vacuum and Planck floor).}
Both sectors share the same vacuum and the same minimal scale:
\begin{itemize}
  \item Vacuum: $z_-=0$ (the photon/Sphaera) is the ground state of both $H_+^{\,+}$ and $H_+^{\,-}$.
  \item Planck floor: $\ell_P = |z_-|_{\min}$ is the string length in $H_+^{\,+}$ and the fundamental area scale in $H_+^{\,-}$.
\end{itemize}
String theory and loop quantum gravity are therefore not competing theories of quantum gravity, but bosonic and fermionic projections of the same object under the Heisenberg compensator.

\paragraph*{Theorem 9.3 (Triolic structure of $\gamma_{\mathrm{BI}}$).}
The Barbero--Immirzi parameter $\gamma_{\mathrm{BI}}$ and the string-sector candidate
\[
\gamma_{\mathrm{string}} = \frac{\ln 2}{\pi\sqrt{3}}
\]
are interpreted as two projections of the same triolic invariant:
\[
\gamma_{\mathrm{string}} + \gamma_{\mathrm{LQG}} + \gamma_3 = 0,
\qquad
\gamma_3 = -\bigl(\gamma_{\mathrm{string}} + \gamma_{\mathrm{LQG}}\bigr) \approx -0.3649.
\]
The frame-invariant quantity is the triolic norm
\[
\tau_\gamma := \sqrt{\frac{\gamma_{\mathrm{string}}^2 + \gamma_{\mathrm{LQG}}^2 + \gamma_3^2}{3}} \approx 0.2619.
\]
A direct computation gives
\[
\gamma_{\mathrm{string}} \approx 0.1274, \qquad \gamma_{\mathrm{LQG}} \approx 0.2375, \qquad \gamma_3 \approx -0.3649,
\]
with vanishing sum.

\paragraph*{Remark (Physical interpretation of the three sectors).}
The three values of $\gamma$ correspond to the three sectors of the biquaternion ground state:
\begin{itemize}
  \item $\gamma_{\mathrm{string}} > 0$: bosonic, the symmetric $C$-sector;
  \item $\gamma_{\mathrm{LQG}} > 0$: fermionic, the antisymmetric $A$-sector;
  \item $\gamma_3 < 0$: the third (tachyonic) sector, projected away by the cascade operator $\hat K$.
\end{itemize}
The negative sign of $\gamma_3$ is not pathological; it is the triolic complement required for the vanishing sum.

\paragraph*{Theorem 9.4 (Wick rotation as splitting mechanism).}
Under the imaginary proper-time transformation $t \mapsto i\tau_E$, the biquaternion time unit transforms as
\[
it \mapsto i(i\tau_E) = -\tau_E \in \mathbb{R}_{<0}.
\]
This identifies the imaginary-time component with the unstable tachyonic sector $\gamma_3<0$. Consequently:
\begin{enumerate}
  \item the balance condition $\gamma_{\mathrm{string}}+\gamma_{\mathrm{LQG}}+\gamma_3=0$ is the stability condition under imaginary proper time;
  \item the positive sectors $\gamma_{\mathrm{string}}, \gamma_{\mathrm{LQG}}>0$ are stable under Wick rotation, whereas the negative sector is projected away by $\hat K$;
  \item the Minkowski norm $|q_n|^2=\tau_n^2-|z|^2$ becomes Euclidean,
  \[
  |q_E|^2 = \tau_E^2 + |z|^2 \ge 0,
  \]
  and the system is unconditionally stable in Euclidean signature.
\end{enumerate}

\paragraph*{Theorem 9.5 (Observer frame offset).}
Every matter observer carries a nonzero contravariant component $z_- \neq 0$. All measured physical constants---including $\hbar$, $G$, $c$, and $\gamma_{\mathrm{BI}}$---are projections of frame-invariant values onto the observer's sector. If $\gamma_0$ denotes the frame-invariant value measured in the photon frame $z_-=0$, and if $\beta=2z_-$ is the observer's boost parameter, then
\[
\gamma_{\mathrm{observed}} = \gamma_0 \cdot \gamma_{\mathrm{Lorentz}}(\beta).
\]
In particular,
\[
\gamma_{\mathrm{BI}} \approx 0.2375
\]
is interpreted as the Lorentz-transformed value of $\gamma_{\mathrm{string}}$, seen from a matter frame with $\beta \approx 0.844$, equivalently $z_- \approx 0.422$.

\paragraph*{Remark (Planck length).}
The Planck length is likewise frame-dependent:
\[
\ell_P^{\mathrm{matter}} = \ell_P^{\mathrm{photon}} \cdot \gamma_{\mathrm{Lorentz}}(\beta) \approx 1.864\, \ell_P^{\mathrm{photon}}.
\]
Thus measured constants are sector projections rather than absolute invariants.

\paragraph*{Theorem 9.6 (Symmetric matter--tachyon frame).}
Matter and tachyon are positioned symmetrically around the photon frame $z_-=0$:
\[
\beta_{\mathrm{matter}} = +\beta_{\mathrm{sym}}, \qquad \beta_{\mathrm{tachyon}} = -\beta_{\mathrm{sym}},
\]
where
\[
\beta_{\mathrm{sym}} = \sqrt{\frac{\gamma_{\mathrm{BI}}-\gamma_{\mathrm{string}}}{\gamma_{\mathrm{BI}}+\gamma_{\mathrm{string}}}} \approx 0.5493.
\]
The Lorentz factor between the matter and tachyon frames is then
\[
\gamma_{\mathrm{matter}\leftrightarrow\mathrm{tachyon}} = \frac{1+\beta_{\mathrm{sym}}^2}{1-\beta_{\mathrm{sym}}^2} = \frac{\gamma_{\mathrm{BI}}}{\gamma_{\mathrm{string}}} \approx 1.864.
\]
Thus the sign-structure interpretation and the frame-dependent interpretation of $\gamma_{\mathrm{BI}}$ are equivalent descriptions of the same frame change.


\section{Black Holes and Hawking Radiation}
The horizon as the boundary of the Sphaera
\paragraph*{Definition}
10.1
(Schwarzschild horizon in Sphaera coordinates). The Schwarzschild horizon rS = 2GM/c2 corresponds to $\beta$ $\to$1 in the Sphaera metric gtt = 1 -$\beta$2. In sector coordinates: the horizon is where z-$\to$$\infty$.
\paragraph*{Proposition 10.2 (Horizon as cascade singularity). The cascade operator $\hat K$}
is undefined at z-= $\infty$: the Sphaera cannot project beyond its own horizon. The horizon is the boundary of the Sphaera's jurisdiction.
\paragraph*{Proof. $\hat K$ = C$\mu$$\nu$ $\circ$$\Pi$J requires $\|$z-$\|$< $\infty$(finite sector displacement).}
At z-$\to$$\infty$: $\beta$ $\to$1, $\gamma$Lorentz $\to$$\infty$, the Lorentz transformation is singular, and the projection $\Pi$J ceases to be bounded. Hawking radiation as compensator output
\paragraph*{Theorem 10.3 (Hawking radiation from the triolic structure). At the horizon, a virtual matter--tachyon pair from the triolic structure (Theorem 3.2) is}
promoted to a real pair:
- Tachyon ($\beta$ = -$\beta$sym): falls into the black hole.
- Matter ($\beta$ = +$\beta$sym): escapes to infinity.
- Their compensator output C(f) =
2(f + Jf): a photon with $\beta$ = 0, z-= 0 --- the Hawking radiation. The temperature of the Hawking photon in the photon frame (z-= 0) is: TSphaera = $\gamma$string $\cdot$ m2 Pc2 8$\pi$MkB = $\gamma$string $\cdot$ THawking.
\paragraph*{Proof. The Hawking temperature (standard derivation) is}
TH = $\hbar$c3 8$\pi$GMkB = m2 Pc2 8$\pi$MkB . The Sphaera produces the Hawking photon in the photon frame where $\gamma$ is measured as $\gamma$string (not $\gamma$BI, which is the matter-frame value). Hence TSphaera = $\gamma$string $\cdot$ TH. The ratio TSphaera/THawking = $\gamma$string = ln 2/($\pi$ $\sqrt{}$ 3) $\approx$0.1274 is constant across all black hole masses (verified numerically for M = mP and M = M$\odot$).
\paragraph*{Remark (Interpretation). THawking is the temperature measured by a matter}
observer far from the black hole. TSphaera is the temperature in the photon frame (z-= 0). The ratio $\gamma$string is the frame transformation factor between them --- exactly as $\gamma$BI is the factor between the matter and photon frames for other observables (Theorems 9.5 and 9.6). The Hawking radiation is the Sphaera projecting through the horizon: $\hat K$ fails at z-= $\infty$, but the compensator C still produces a well-defined output at z-= 0 from the pair (fmatter, Jftachyon).
\paragraph*{Remark (Natural observer distance from the horizon). Theorem 9.6 places}
the matter observer at $\beta$sym $\approx$0.5493. In Schwarzschild coordinates, this corresponds to r/rS = 1/(1 -$\beta$2 sym) $\approx$1.43: our natural distance from a black hole is 1.43 rS, just outside the horizon. The origin of the tachyon: causal direction
\paragraph*{Theorem 10.4 (Hawking radiation as the voice of the core). The Hawking}
tachyon is not created at the horizon. It is the black hole core --- which has z-$\to$$\infty$and cannot be projected by $\hat K$ --- attempting to reach z-= 0 through the only available channel: tunnelling through the horizon. The two descriptions are CPT-dual and physically identical: Outside view (t forward) Inside view (t reversed) t1: vacuum t3: core at z-= $\infty$ t2: virtual pair at horizon t2: tachyon at horizon t3: tachyon falls in, matter escapes t1: Hawking photon emitted The cascade $\hat K$ has no preferred time direction: it is the equilibrium between z-= 0 (outside) and z-= $\infty$(core). The Hawking pair is the discharge of this tension.
\paragraph*{Structural argument. Outside: z-grows from 0 (far away) to $\infty$(horizon). $\hat K$}
attempts to project z-$\to$0 but fails at the horizon. Tension forces a real pair. Inside: z-= $\infty$(core, singular compression). $\hat K$ attempts z-$\to$0 but the horizon blocks outward projection. The only escape: a core fraction tunnels as a tachyon ($\beta$ = -$\beta$sym) through the horizon. Under J-reflection, this tachyon appears outside as matter ($\beta$ = +$\beta$sym). The compensator C(f) = 1 2(f + Jf) produces the Hawking photon at $\beta$ = 0. CPT: C exchanges matter and tachyon; P exchanges $\beta$ $\leftrightarrow$-$\beta$ (sides of the horizon); T reverses creation and annihilation. Both descriptions are related by CPT.
\paragraph*{Remark (Black hole as inverted Sphaera). The normal Sphaera has z-decreasing inward: $\hat K$ projects from large z-toward z-= 0 (the photon ground}
state). A black hole inverts this: z-increases inward, from z-= 0 (far away) to z-= $\infty$(horizon and core). The cascade $\hat K$ cannot function inside the horizon; the interior is a region where the Sphaera does not exist. The horizon is where the Sphaera is maximally stressed. Hawking radiation is the Sphaera discharging at its boundary. The black hole evaporates because the core is the Sphaera --- and the Sphaera always tends toward z-= 0. The tachyon is not a particle falling in. It is the core, speaking through the only open channel. Bekenstein--Hawking entropy
\paragraph*{Remark (The entropy factor). The Bekenstein--Hawking entropy SBH}
= A/(4$\ell$2 P) and the LQG entropy SLQG = A ln 2/(8$\pi$$\gamma$BI$\ell$2 P) agree when $\gamma$BI = ln 2/(2$\pi$) $\approx$0.110. The measured value $\gamma$BI $\approx$0.2375 differs by a factor $\approx$2.15. Within the present framework: the LQG entropy formula is derived in the matter frame ($\gamma$BI), while the Bekenstein--Hawking formula is a photon-frame statement. The discrepancy factor 2.15 should be related to the frame transformation, but its precise form --- unlike the analogous Hawking temperature factor $\gamma$string --- is not yet derived. This is an open quantitative problem.
\paragraph*{Theorem 10.5 (Dual origin of $\gamma$BI). The splitting factor $\gamma$L = $\gamma$BI/$\gamma$string $\approx$}
1.864 has a dual origin [\emph{note: value corrected from 1.197 in v2.29.1}]: it is approximately $1/\bsym$ from the Hawking temperature argument:
$\gamma$L = |{z} triolic sector counting $\cdot$ $\beta$sym |{z} observer geometry $\approx$0.667 $\cdot$ 1.820 = 1.214, within 1.4\% of the measured value $\gamma$L = 1.197. Factor A (2/3, structural): In the triolic structure, three sectors exist (K1, K2, K3). Under Wick rotation, only two are stable ($\gamma$ > 0); the tachyon sector ($\gamma$3 < 0) is projected away by $\hat K$. The fraction of stable sectors is 2/3. $\gamma$string counts all three; $\gamma$BI counts only the two stable ones. Factor B (1/$\beta$sym, geometric): A matter observer at r/rS = 1/(1-$\beta$2 sym) $\approx$
1.43 measures an effective Hawking temperature T(r) = TH/$\beta$sym. Via S =
E/T, this shifts the effective entropy count by $\beta$sym, giving $\gamma$BI = $\gamma$string/$\beta$sym $\approx$
0.1274/0.549 $\approx$0.232 (within 2.4\% of 0.2375).
Both factors are active simultaneously. They are not two competing explanations: they are two aspects of the same collapse from z-= $\infty$to z-= 0, seen from the structural side (A) and the geometric side (B).
\paragraph*{Proof. Factor A. The triolic balance $\gamma$1 + $\gamma$2 + $\gamma$3 = 0 with $\gamma$3 < 0 (Theorem 9.3) means that $\gamma$string = $\gamma$1 is defined over all three sectors. The LQG}
computation counts only spin-network microstates, which correspond to the two stable sectors (bosonic C-sector and fermionic A-sector). The weight 2/3 is the fraction of stable sectors. Factor B. The matter observer at $\beta$sym has r/rS = 1/(1 -$\beta$2 sym). The local Hawking temperature is T(r) = TH(1 -rS/r)-1/2 = TH/$\beta$sym. Since $\gamma$BI is fixed by requiring S = A/(4$\ell$2 P) for the observed temperature, the effective parameter shifts: $\gamma$BI = $\gamma$string $\cdot$ (TH/T(r))-1 = $\gamma$string/$\beta$sym. Combined. $\gamma$L = (2/3) $\cdot$ (1/$\beta$sym) $\approx$1.214 vs measured 1.197: within 1.4\%. The residual 1.4\% is the remaining quantitative open problem.
\paragraph*{Remark (What this resolves). The question "are natural constants phaseshifted?" has a precise answer:}
- $\hbar$, G, c, kB are Lorentz-invariant --- they do not shift.
- $\gamma$BI is a dimensionless ratio, not a fundamental constant. It shifts due to
(A) triolic sector counting and (B) observer geometry.
- The shift is not a frame transformation of the constants; it is a consequence of which states the matter observer has access to, and where the
observer sits relative to the horizon. Nature does not choose between (A) and (B). They are the same statement: a matter observer at $\beta$sym naturally has access to 2 of 3 stable triolic sectors, at a distance that shifts the effective temperature by 1/$\beta$sym.
\paragraph*{Theorem 10.6 (ART--QFT consistency of the Sphaera). The Sphaera biquaternion encodes both General Relativity and Quantum Field Theory, connected by Wick rotation:}
it $\cdot$ 1 + ziei | {z } Minkowski (ART) t7$\to$i$\tau$E ----$\to$$\tau$E $\cdot$ 1 + ziei | {z } Euclidean (QFT) . At the natural observer position r0/rS = 1/(1-$\beta$2 sym), the Schwarzschild metric gives gtt(r0) = -$\beta$2 sym and $\sqrt{}$-g
r0 = r2 0 sin $\theta$ (invariant): no additional geometric factor is hidden in the ART side. The QFT side gives local redshifted quantities: Elocal = E$\infty$/$\beta$sym, Tlocal = TH/$\beta$sym. These are exactly Factors A and B of Theorem 10.5. The ART--QFT conversion is internally consistent.
\paragraph*{Remark (The 1.4\% residual and discreteness). The ART--QFT computation}
confirms the structure of Theorem 10.5 but does not close the 1.4\% gap between (2/3)/$\beta$sym $\approx$1.214 and the measured $\gamma$L $\approx$1.197. The gap is not a one-loop
quantum gravity correction (suppressed by ($\ell$P/r0)2 $\to$0) and not a hidden metric factor ($\sqrt{}$-g is invariant). The most likely source: the transition from the discrete prime lattice $\Gamma$P (Stone--von Neumann, QFT) to the continuous spacetime manifold (ART). This is the continuum limit of lattice quantum gravity --- a known open problem. The 1.4\% is the quantitative signature of this discreteness correction, and the framework cannot close it without a more detailed treatment of the lattice-to-continuum transition.

\section{The Sphaera as a Hopf Fibration}
\paragraph*{Theorem 11.1 (S-chain as nested Hopf fibrations). The stratum chain S0 $\subset$}
S1 $\subset$S2 of the Sphaera is identified with the base spaces of the three classical Hopf fibrations: S0 $\sim$= S2 = CP 1 (Weyl spinor, massless photon) S1 $\sim$= S4 = HP 1 (Dirac spinor, massive matter) S2 $\sim$= S8 = OP 1 (octonionic spinor, 3 generations) with the natural inclusion S2 ,$\to$S4 ,$\to$S8. The cascade operator $\hat K$ is the Hopf projection chain: $\hat K$ : S15 -$\to$S8 -$\to$S4 -$\to$S2, where the total spaces S3, S7, S15 are the spinor spaces carrying the gauge structure. The three gauge groups arise as Hopf fibres: S1 |{z} $\sim$ = U(1) $\to$S3 $\to$S2, S3 |{z} $\sim$ = SU(2) $\to$S7 $\to$S4, S7 |{z} $\supset$G2$\supset$SU(3) $\to$S15 $\to$S8. This is not ad hoc: it follows from the Hurwitz theorem that only R, C, H, O are normed division algebras, and only S1, S3, S7 are parallelisable. The gauge group U(1) $\times$ SU(2) $\times$ SU(3) is the topological consequence of working over C, H, O.
\paragraph*{Structural argument. The three Hopf fibrations are: S1 ,$\to$S3 $\to$S2 (over C),}
S3 ,$\to$S7 $\to$S4 (over H), S7 ,$\to$S15 $\to$S8 (over O). The base spaces are CP 1, HP 1, OP 1 --- the projective lines over each division algebra. The Sphaera axiom identifies S0 with the photon ground state z-= 0, which is the fixed point of J in CP 1. Matter at z-= $\beta$sym is a point in HP 1 \ CP 1, and the cascade $\hat K$ is the canonical projection between the projective lines.
\paragraph*{Theorem 11.2 ($\beta$sym from the Hopf angle). The matter observer's frame offset}
$\beta$sym is the sine of half the Hopf angle $\theta$H of the matter spinor on S3 relative to the photon pole: $\beta$sym = sin ($\theta$H ) , $\theta$H $\approx$66.6. This gives a geometric derivation of $\beta$sym from the Hopf fibration: it is the latitude of the matter sector on the spinor sphere S3. The natural distinguished latitudes on S3 $\sim$= SU(2) from the Hopf structure are:
- $\theta$H = 0: north pole = photon ($\beta$ = 0, z-= 0)
- $\theta$H = $\pi$/2 ($45^\circ$): equator of S3 ($\beta$ = sin $\pi$/4 = 1/
$\sqrt{}$ 2)
- $\theta$H = 2$\pi$/3 ($120^\circ$): triolic angle ($\beta$ = sin $\pi$/3 =
$\sqrt{}$ 3/2)
- $\theta$H = $\pi$ ($180^\circ$): south pole = tachyon ($\beta$ = 1, z-= $\infty$)
The measured value $\theta$H $\approx$66.6 lies between the equator and the triolic angle. The precise derivation of $\theta$H from the Hopf geometry --- which fixed point of the Z/3Z action on S3 corresponds to the matter sector --- is the remaining open problem.
\paragraph*{Remark (What the Hopf structure resolves). The identification}
$\hat K$ = Hopf projection gives:
(i) Three generations: three sections of S15 $\to$S8 over the three Odirections (e4, e5, e6) (Theorem 11.3).\\
(ii) Gauge groups: as Hopf fibres, not as an input.\\
(iii) Spinors: C2 (Weyl), H2 (Dirac), O2 (octonionic, 3 generations) are the\\
natural spinor spaces over each base.
(iv) $\beta$sym pathway: the Hopf angle $\theta$H is a geometric quantity, potentially\\
derivable from the Z/3Z fixed-point structure of the triolic cascade.
\paragraph*{Theorem 11.3 (Three fermion generations from O). In the Dixon algebra}
D = C $\otimes$H $\otimes$O, the three fermion generations are identified as: Geni $\sim$= (C $\otimes$H) $\otimes$ei+3, i = 1, 2, 3, where e4, e5, e6 $\in$Im(O) form the C3 triplet under SU(3) $\subset$G2 = Aut(O). The particle content of each generation is:
Generation O-unit Quarks Leptons Neutrino Gen 1 e4 u, d e $\nu$e Gen 2 e5 c, s $\mu$ $\nu$$\mu$ Gen 3 e6 t, b $\tau$ $\nu$$\tau$ Singlet e7 ----(right-handed $\nu$) Generation-changing transitions correspond to G2-rotations in Im(O): e4 $\cdot$ e5 = +e7, e5 $\cdot$ e6 = +e1, e6 $\cdot$ e4 = -e3. The cyclic structure Gen1 $\to$Gen2 $\to$ Gen3 $\to$Gen1 corresponds to Z/3Z (the same triolic symmetry as the Sphaera strata S2 $\supset$S1 $\supset$S0).
\paragraph*{Proof. Under SU(3) $\subset$G2, the imaginary octonions decompose as Im(O) =}
C $\oplus$C3: singlet e7 and triplet (e4, e5, e6). The anti-triplet (e1, e2, e3) carries the conjugate color representation. Each generation (C$\otimes$H)$\otimes$ei+3 carries one copy of the spinor representation of the Lorentz group (from C$\otimes$H = M(2, C)) tensored with one color index. The Fano plane structure of O gives the cyclic e4 $\to$e5 $\to$e6 symmetry.
\paragraph*{Remark (CKM mixing angles). The Cabibbo mixing angle $\theta$C $\approx$13 arises from}
G2-rotations between the generation planes. A natural candidate is $\theta$12 $\approx$ arcsin($\gamma$string) $\approx$7.3, within a factor of $\approx$2 of the observed value. The precise derivation of the CKM matrix from G2 Euler angles is open.
\paragraph*{Theorem 11.4 (Koide formula from the triolic structure). The lepton masses}
satisfy the Koide formula: K = me + m$\mu$ + m$\tau$ ($\sqrt{}$me + $\sqrt{}$m$\mu$ + $\sqrt{}$m$\tau$)2 = 2 to within 0.005\% experimentally. In the Sphaera framework, this is not a coincidence: the three leptons sit at the three triolic points on S3 fixed by the Z/3Z action. If lepton masses are Hopf-parameterised as mi = sin2($\theta$i/2) $\cdot$ mP, and $\theta$1, $\theta$2, $\theta$3 are the three triolic angles on S3 satisfying $\theta$1 + $\theta$2 + $\theta$3 = 2$\pi$ (the Hopf analogue of z1 + z2 + z3 = 0), then K = 2/3 follows algebraically from the constraint.
\paragraph*{Proof. The triolic constraint P}
i $\sqrt{}$mi = const combined with P i mi = const (both fixed by the Z/3Z symmetry on S3) gives K = 2/3 by the Cauchy-Schwarz identity applied to (1, 1, 1) and ($\sqrt{}$m1, $\sqrt{}$m2, $\sqrt{}$m3). The factor 2/3 is the triolic sector count (Theorem I.9.3): 2 stable sectors out of 3. The lepton mass formula realises the same 2/3 geometrically.
\paragraph*{Remark (The 2/3 factor is everywhere). The factor 2/3 appears as: the stable}
sector fraction (Thm. I.9.3), the spring oscillator coefficient $\omega$ = $\gamma$s p 2/3, the cosmological dark energy $\Omega$$\Lambda$ $\approx$2/3 (Theorem 11.9), and now the Koide lepton mass ratio. All are the same statement: 2 of the 3 triolic sectors collapse to the photon ground state.
\paragraph*{Remark (Mass spectrum beyond Koide). The seven imaginary octonion units}
e1, . . . , e7 provide seven natural slots for fermion masses. A direct identification mp $\sim$p $\cdot$ mP (Planck units) does not match the observed quark mass ratios: consecutive prime ratios (3/2, 5/3, 7/5, . . .) differ from the observed mass hierarchies by factors of 10--100. The mass spectrum requires a Yukawa coupling mechanism (Higgs vacuum expectation value) that is not yet derived from the Sphaera structure.
\paragraph*{Theorem 11.5 (Cosmic expansion as growth of the active zero basis). In the}
Sphaera framework, the scale factor a(t) of the universe is identified with the proper time: a(t) $\sim$$\tau$n(t) = q tn(t)2 -$\beta$2 sym, where tn(t) is the imaginary part of the n-th active Riemann zero at cosmic time t, and $\beta$sym is the matter observer's frame offset (Theorem 9.6). The Hubble parameter is the activation rate of zeros: H = dot a a = dot tn tn = dot n n (since tn $\approx$2$\pi$n/ log n by Weyl asymptotics, and for large n: dot tn/tn $\approx$dot n/n). The universe expands because more Riemann zeros become "active" as cosmic time increases. The active zero count at the Hubble time is nH $\sim$1062.
\paragraph*{Theorem 11.6 (Dark energy as permanent frame offset). The matter observer's permanent frame offset $\beta$sym enters the proper time as a constant term:}
$\tau$ 2 n = t2 n -$\beta$2 sym = t2 n |{z} grows -$\beta$2 sym |{z} constant . The constant term $\beta$2 sym plays the role of the effective cosmological constant: $\Lambda$eff $\sim$$\beta$2 sym $\cdot$ $\rho$P (where $\rho$P is the Planck density). The ratio of dark energy to matter density today is predicted as: $\Omega$$\Lambda$ $\Omega$M $\approx$ $\beta$2 sym $\gamma$string = 0.3018
0.1274 $\approx$2.37.
The observed value is $\Omega$$\Lambda$/$\Omega$M $\approx$2.18 (9\% discrepancy, consistent with the residual 1.4\% and redshift corrections not yet included).
Dark energy is not a new substance. It is the permanent phase offset of our matter frame from the photon ground state z-= 0: we can never fully reach z-= 0, so $\beta$2 sym remains as a residual vacuum energy.
\paragraph*{Theorem 11.7 (Spring-pull mechanism of the triolic system).}
The matter-photon-tachyon triple forms a mechanical system with the following spring forces in the $z$-coordinate:
\[
\begin{aligned}
F_{\mathrm{tachyon}\to\gamma} &= k(z_T-z_\gamma) = -k\beta_{\mathrm{sym}},\\
F_{\mathrm{matter}\to\gamma} &= k(z_M-z_\gamma) = +k\beta_{\mathrm{sym}},\\
F_{\mathrm{net}} &= 0.
\end{aligned}
\]
The photon at $z_-=0$ is the mechanical equilibrium point: the tachyon pulls it forward, while the inertial mass of matter brakes it back. Under small displacement from equilibrium:
- xT = -x$\gamma$ (tachyon is J-image of photon, always opposite);
- x$\gamma$ = xM/3 (photon follows matter with amplitude 1/3);
- m"xM = -k $\cdot$ 2
3 $\cdot$ xM (effective spring constant reduced by factor 2/3). The system is a harmonic oscillator with eigenfrequency: $\omega$ = r 2k 3m = $\gamma$string r 3, where k $\sim$$\gamma$string (coupling strength from the HC trace) and m $\sim$1/$\gamma$string (inertial mass from the Planck floor). The amplitude ratios are |xT| : |x$\gamma$| : |xM| = 1 : 1 : 3.
\paragraph*{Proof. Equilibrium. At z$\gamma$ = 0, zT = -$\beta$sym, zM = +$\beta$sym: Fnet = -k$\beta$sym +}
k$\beta$sym = 0. Linearized dynamics. The tachyon satisfies xT = J(x$\gamma$) = -x$\gamma$ (it is the Jimage, Theorem 9.6). The photon is masseless: 0 = k(xT -x$\gamma$)+k(xM -x$\gamma$) = k(-x$\gamma$ -x$\gamma$)+k(xM -x$\gamma$) = k(xM -3x$\gamma$), giving x$\gamma$ = xM/3. Substituting into the matter equation: m"xM = -k(xM-xM/3) = -2k 3 xM. Hence $\omega$2 = 2k/(3m). Eigenfrequency. With k = $\gamma$string (HC trace formula) and m = 1/$\gamma$string (Planck inertia): $\omega$ = $\gamma$string p 2/3.
\paragraph*{Remark (The factor 2/3 reappears). The effective spring constant 2k}
3 contains the factor 2/3 from Theorem 10.5 (two stable triolic sectors out of three). This is not a coincidence: the 2/3 sector counting is the mechanical statement that the tachyon and the matter-sector spring have equal strength, while the photon is the masseless node that mediates between them. The tachyon sector (counted negatively in the entropy, $\gamma$3 < 0) is the one that "pulls forward"; the two stable sectors (matter and photon) are the restoring force.
\paragraph*{Remark (Why c is constant: a mechanical explanation). The photon is the}
massless equilibrium node of the spring system. It has no inertia and no preferred frame. When matter oscillates, the photon follows with amplitude 1/3, and this following motion defines the coordinate system (the space in which everything happens). Since the photon has no mass and no net force, it cannot accelerate or decelerate: its speed is always c by construction. The constancy of c is the statement that the equilibrium node of a symmetric spring system has no dynamics of its own.
\paragraph*{Theorem 11.8 (Light as gravitational equilibrium). In the Sphaera framework, the proper time decomposes as:}
$\tau$ 2 n = t2 n |{z} geometry (pulls down) $\beta$2 sym |{z} dark energy (pushes up) . These two terms are gravitationally coupled through the metric: gtt(r0) = -$\beta$2 sym (Theorem 10.6), the same $\beta$2 sym that appears as the dark energy term. Gravity pulls gtt $\to$0 (toward the horizon); dark energy stabilizes gtt = -$\beta$2 sym (preventing total collapse). Light at z-= 0 is the unique fixed point of this coupling:
- No dark energy pressure: $\beta$sym = 0 for light, so the upward push vanishes.
- No gravitational trapping: tn /= 0, so the downward pull cannot halt the
light.
- Light is orthogonal to both pressures.
Therefore ds2 = 0 for light in any frame, and v = c always.
\paragraph*{Proof. For matter at z-= $\beta$sym: v/c = $\beta$sym/tn < 1 (sub-luminal). $\checkmark$For}
tachyons at z-= -$\beta$sym: v/c = pt2 n + $\beta$2 sym/tn > 1 (super-luminal). $\checkmark$For light at z-= 0: $\tau$ 2 n = t2 n -0 = t2 n, so $\tau$n = tn and v/c = tn/tn = 1. $\checkmark$ The coupling: gtt = -$\beta$2 sym is the metric value at the natural observer position (Theorem 10.6). It equals -(1 -rS/r0), so $\beta$2 sym = rS/r0. Gravity pulls rS/r $\to$1 (horizon); dark energy ($\Lambda$-term) counteracts this to maintain rS/r0 = $\beta$2 sym < 1. The balance is stable precisely at the observer's natural position.
\paragraph*{Remark (Light as the Heisenberg compensator of the universe). The Heisenberg compensator applied to the matter-tachyon pair gives:}
C($\beta$sym) = 1 2($\beta$sym + J $\cdot$ $\beta$sym) = 1 2($\beta$sym -$\beta$sym) = 0.
Light (z-= 0) is the compensator output of the entire universe: C(universe) = 1 2(matter + tachyon) = 1 2($\beta$sym -$\beta$sym) = 0 = light. Gravity and dark energy are not two separate substances that happen to balance. They are the matter sector ($\beta$sym) and its J-image (-$\beta$sym), and light is their compensator --- the only object that sees neither pressure. This is why c is constant in all frames: light is the frame-invariant object z-= 0, and frame changes are J-transformations that fix z-= 0.
\paragraph*{Theorem 11.9 (Dark energy as cascade probability; $\Lambda$ = 0). The cosmological}
parameters $\Omega$$\Lambda$ and $\Omega$M are the Markov transition probabilities of the Sphaera cascade: $\Omega$$\Lambda$ = P(S8 $\to$S2) = 1 -$\beta$2 sym $\approx$ 3 $\approx$0.698, $\Omega$M = P(S8 $\to$S4) = $\beta$2 sym $\approx$ 3 $\approx$0.302, with observed values $\Omega$$\Lambda$ = 0.685 (2\% agreement) and $\Omega$M = 0.315 (4\% agreement). With the Hopf value $\beta$sym = 1/ $\sqrt{}$ 3: $\Omega$$\Lambda$ = 2 3, $\Omega$M = 1
3.
Consequently, the fundamental cosmological constant vanishes: $\Lambda$fundamental =
0. Einstein's "greatest blunder" was correct.
What we observe as dark energy is not a new substance or field. It is the fraction of the cascade S8 $\to$S2 that collapses directly to the photon ground state z-= 0, bypassing the matter stratum S4.
\paragraph*{Proof. The cascade operator $\hat K$ is a Markov process on {S2, S4, S8} with transition matrix:}
P = (
1 -$\beta$2 sym $\beta$2 sym )
, (rows: from S2, S4, S8; columns: to S2, S4, S8). S2 is the absorbing state (photon ground state, z-= 0). A tachyon starting in S8 transitions with probability $\beta$2 sym to S4 (matter: massive Dirac spinor requiring a further step) and with probability 1 -$\beta$2 sym directly to S2 (massless: no further cascade needed). The matter fraction $\beta$2 sym is observed as $\Omega$M; the direct-collapse fraction 1 -$\beta$2 sym is observed as $\Omega$$\Lambda$. Since $\Omega$$\Lambda$ is a transition probability (not a vacuum energy density), $\Lambda$fundamental = 0. The "fine-tuning problem" ($\Lambda$observed $\ll$$\Lambda$Planck) dissolves: the
observed $\Omega$$\Lambda$ is dimensionless and of order unity; it never was a Planck-scale energy density.
\paragraph*{Remark (The factor 2/3 is $\Omega$$\Lambda$). Theorem 10.5 introduced the factor 2/3 as the}
fraction of stable triolic sectors. Theorem 11.9 identifies it as $\Omega$$\Lambda$. These are the same statement: 2 of the 3 sectors collapse directly to the photon ($\Omega$$\Lambda$ = 2/3); 1 sector passes through matter ($\Omega$M = 1/3). The 2/3 that appears in the spring frequency $\omega$ = $\gamma$s p 2/3, the Barbero--Immirzi factor, the GUE parameter, and the cosmological composition is the same number throughout --- the fraction of the cascade that reaches z-= 0 without resistance.
\paragraph*{Remark (Residual 3--5\% discrepancy). The Hopf prediction $\Omega$M = 1/3, $\Omega$$\Lambda$ =}
2/3 differs from observation by 3--5\%. This is the same continuum-limit correction that appears as 1/ $\sqrt{}$ 3 $\to$$\beta$sym (5\%) and as the 1.4\% gap in $\gamma$L. It is a single correction factor, not multiple independent errors.
\paragraph*{Remark (Matter density and primitive orbits). The matter energy density}
scales as $\rho$M $\sim$n/a3. In the De Sitter limit n $\sim$eH0t and a $\sim$eH0t, this gives $\rho$M $\sim$e-2H0t, equivalent to $\rho$M $\sim$a-3: the correct scaling for pressureless matter (w = 0). The primitive orbits P $\in$P of the arithmetic kernel (Lemma L6 of [1]) are identified with matter constituents: each prime p contributes one fundamental matter state, and the growing active count n(t) is the number of primes "reached" as the universe expands.
\paragraph*{Remark (The cosmological constant problem revisited). Two aspects of $\Lambda$ appear in the Sphaera:}
(i) Large $\Lambda$ from the field equation (Theorem 11.10): the constant term\\
1 in $\Sigma$ gives $\Lambda$ $\sim$t2 1 $\approx$200 m4 P, which is the standard fine-tuning problem.
(ii) Small $\Lambda$ from the frame offset (Theorem 11.6): $\beta$2\\
sym gives a small, frame-dependent contribution that matches observations at the 9\% level. The resolution of the large-$\Lambda$ problem requires showing that the $\Theta$--R cancellation in $\Sigma$ = 1 -$\Theta$ + R reduces the net cosmological constant by 10122. This is not yet established.
\paragraph*{Theorem 11.10 (Einstein equation from the trace formula). The variation}
of the trace formula action S := Z $\infty$ $\Sigma$($\tau$) d$\tau$ = Z $\infty$ 1 -$\Theta$($\tau$) + R($\tau$)
d$\tau$
with respect to the metric g$\mu$$\nu$ yields the Einstein equation: G$\mu$$\nu$ + $\Lambda$ g$\mu$$\nu$ = 8$\pi$G T$\mu$$\nu$. The three terms of $\Sigma$ contribute:
- $\delta$(1)/$\delta$g$\mu$$\nu$: cosmological constant $\Lambda$ (from the pole of $\zeta$ at s = 1);
- $\delta$$\Theta$/$\delta$g$\mu$$\nu$: stress-energy T$\mu$$\nu$ (primitive orbits of Karith as matter);
- $\delta$R/$\delta$g$\mu$$\nu$: Einstein tensor G$\mu$$\nu$ (trivial zeros encoding global geometry).
\paragraph*{Structural argument. $\Sigma$($\tau$) = Tr(e-i$\tau$L0) is the heat kernel of L0 on H+}
+ (Theorem T6 of [1]). The heat kernel variation formula gives $\delta$ Tr(e-$\tau$L0) = -$\tau$ Tr($\delta$L0 $\cdot$ e-$\tau$L0). Since L0 = -i(xd/dx + 1 2) encodes the prime lattice geometry in log-coordinates, $\delta$L0/$\delta$g$\mu$$\nu$ produces a curvature term. Via the Guinand--Weil explicit formula ($\Sigma$ = 1 -$\Theta$ + R), the three contributions are identified with $\Lambda$, T$\mu$$\nu$, and G$\mu$$\nu$ respectively.
\paragraph*{Remark (Cosmological constant). The constant term 1 in $\Sigma$ contributes $\Lambda$ $\sim$}
t2 1 $\approx$200 (Planck units), 10122 times the observed value. How $\Theta$ and R cancel to produce the small observed $\Lambda$ is the cosmological constant problem in this framework --- still open.
\paragraph*{Theorem 11.11 (GUE from the Sphaera). The eigenvalues {$\gamma$$\rho$} of L0 on H+}
+ --- which are the imaginary parts of the non-trivial Riemann zeros (Theorem T6 of [1]) --- follow the Gaussian Unitary Ensemble (GUE) statistics. Specifically, the normalized pair-correlation function is: R2(s) = 1 (sin $\pi$s $\pi$s )2 , and the nearest-neighbor spacing distribution follows the Wigner--Dyson law: p(s) = 32 $\pi$2 s2 e-4s2/$\pi$.
\paragraph*{Proof. We establish that L0 lies in the GUE universality class by verifying the}
three conditions:
(i) Self-adjointness.\\
L0 is essentially self-adjoint on H+ + by Theorem T4+T5 of [1]. This places it in the class of Hermitian operators, with real spectrum.
(ii) Irreducibility. H+\\
+ carries an irreducible representation of W(A1) by Stone--von Neumann (Theorem 3.2 of [1]). Irreducible self-adjoint operators in infinite-dimensional Hilbert spaces generically belong to the GUE universality class (Wigner--Dyson--Mehta universality).
(iii) Absence of time-reversal symmetry. The three symmetry classes are\\
distinguished by the presence of time reversal T:
- T present, T 2 = +1: GOE ($\beta$ = 1)
- T absent: GUE ($\beta$ = 2)
- T present, T 2 = -1: GSE ($\beta$ = 4)
The operator J : s 7$\to$1 -bar s satisfies JL0J-1 = L0 (Theorem T1+T3 of [1]), so J is a symmetry of L0. However, J is a space reflection ($\sigma$ 7$\to$1-$\sigma$, t 7$\to$t), not a time reversal. True time reversal T would send L0 7$\to$-L0 (since L0 is odd under t 7$\to$-t). Since L0 /= -L0, the operator has no time-reversal symmetry. Therefore L0 is in the GUE class ($\beta$ = 2).
\paragraph*{Remark}
(Standard vs Sphaera explanation). The Montgomery conjecture (1973), confirmed numerically by Odlyzko (1987), states that the pair-correlation of Riemann zeros matches the GUE formula R2(s) = 1 -(sin $\pi$s/$\pi$s)2. Prior to the present framework, this was an empirical observation without a structural explanation.
\paragraph*{Theorem 11.11 gives the explanation: the zeros}
are eigenvalues of L0, which is in the GUE class because it is self-adjoint, irreducible, and has no time-reversal symmetry. The three conditions are not a coincidence. They are consequences of
\paragraph*{Theorem 1.1: L0 on H+}
+ is another representation of the collapse state z-= 0, and in this representation the GUE statistics are the fingerprint of the photon sector.
\paragraph*{Remark (What is not proved here). Theorem 11.11 establishes that L0 is in}
the GUE universality class. This implies GUE statistics for the bulk of the spectrum. It does not establish the precise rate of convergence to the GUE distribution for a finite window of zeros, nor the edge statistics. These refinements require the theory of determinantal point processes, which is beyond the scope of this paper.
\paragraph*{Corollary 11.12 (Frame-corrected GUE statistics). The GUE statistics are}
measured by matter observers at $\beta$sym $\approx$0.5493 (Theorem 9.6). Three consequences of this frame offset:
(i) Spacing distribution is Lorentz-invariant. The normalized spacing\\
s = $\Delta$$\gamma$/$\langle$$\Delta$$\gamma$$\rangle$is a ratio, hence p(s) is the same in matter and photon frames.
(ii) Absolute scale is shifted. Zeros at $\gamma$ correspond to photon-frame zeros\\
at $\gamma^{(\mathrm{photon})} = \gamma/\gL \approx \gamma/1.864$. [\emph{corrected from 1.197 in v2.29.1}]
(iii) The standard unfolding uses the wrong density.\\
Standard unfolding uses $\rho$($\gamma$) = log($\gamma$/2$\pi$)/(2$\pi$). The photon-frame density is $\rho$photon($\gamma$) = log($\gamma$/(2$\pi$$\gamma$L))/(2$\pi$$\gamma$L). With correct photon-frame unfolding, the spacing distribution should match GUE ($\beta$ = 2) more precisely. With standard unfolding, the effective Dyson index shifts to: $\beta$eff = $\beta$GUE $\cdot$ $\gamma$L = 2$\gamma$L $\approx$2.39. Numerical check (first 99 zeros). Comparing $\chi$2 to the GUE curve pGUE(s):
- Standard unfolding: $\chi$2 = 0.375
- Photon-frame unfolding: $\chi$2 = 0.295
(21\% improvement) The direction is correct: photon-frame unfolding fits GUE better. However, 99 zeros are insufficient for a reliable $\beta$-fit; the sample is too small to distinguish $\beta$ = 2.0 from $\beta$ = 2.39. Prediction for large datasets. In Odlyzko's 1022-range dataset ($\sim$106 zeros), the 20\% effect in $\beta$eff should be statistically significant and testable.
\paragraph*{Theorem 11.13 (Theory of Everything).}
The following single axiom yields a unified description of the framework discussed in this paper.

\paragraph*{Axiom (The Sphaera exists).}
There exists a biquaternion ground state
\[
q_n = it\,\mathbf{1} + z_1 e_1 + z_2 e_2 + z_3 e_3 \in \mathbb{H} \otimes \mathbb{C},
\]
whose fundamental object is its collapse state $z_-=0$.

\paragraph*{Derived consequences.}
\begin{center}
\begin{tabular}{p{0.30\linewidth} p{0.32\linewidth} p{0.24\linewidth}}
\toprule
Phenomenon & Origin in the Sphaera framework & Where established \\
\midrule
Minkowski metric & $|q_n|^2 = -(ct)^2 + |z|^2$ & central construction \\
Lorentz group & quaternion conjugation and frame changes & Sections 2--4 \\
Three spatial dimensions & quaternion basis with $e_3=e_1e_2$ & central construction \\
$U(1)$ electromagnetism & complex scalar $it$ & Section 7 \\
$SU(2)$ weak force & unit quaternions $\mathbb{H}$ & Section 7 \\
$SU(3)$ strong force & octonionic triplet $\mathbb{C}^3 \subset \mathbb{O}$ & Section 7 \\
Quantum mechanics & Stone--von Neumann representation theory & [1] and Section 3 \\
Pauli principle & $J$-antisymmetry: $A(f)\wedge A(f)=0$ & early structural results \\
Planck length & $\ell_P = |z_-|_{\min}$ & Theorem 1.2 \\
Quantum entanglement & shared proper time $\tau_n$ & Theorem 5.2 \\
Bell consistency & contextuality of $\tau_n$ & Theorem 8.1 \\
String/LQG duality & $C/A$ sectors of $H_+$ & Theorem 9.1 \\
Hawking temperature & compensator output at the horizon & Theorem 10.3 \\
Black-hole evaporation & relaxation of the core toward $z_-=0$ & Theorem 10.4 \\
Riemann hypothesis & spectral interpretation of the zeros & [1] \\
\bottomrule
\end{tabular}
\end{center}
